% =====================================================================
%  CUERPO DEL MANUAL DEL OPERARIO - PharmaWeigh
%  Capturas reales del sistema. Sin credenciales.
% =====================================================================

\seccion{1. Acceso al sistema}

\tablaAcceso

\vspace{8pt}

\reglasAcceso

\captura{01_login.png}{Pantalla de entrada. No muestra usuarios ni contraseñas de ejemplo.}

\captura{02_login_lleno.png}{Formulario listo para pulsar \boton{Ingresar}. La contraseña siempre sale enmascarada con puntos: nadie que vea tu pantalla puede leerla.}

\captura{19_login_error.png}{Cuando el correo o la contraseña no coinciden aparece \cd{Credenciales inválidas} sobre el formulario. A los 5 fallos ese correo queda bloqueado 15 minutos.}

\begin{cajaOjo}
\textbf{Tu cuenta es tuya.} Todo lo que se pese, se escanee o se firme con tu usuario
queda registrado con tu nombre, la hora y la IP de la estación. Nunca prestes tu
contraseña, ni siquiera a tu supervisor: él tiene la suya.
\end{cajaOjo}

\avisoValidacion

% =====================================================================
\newpage
\seccion{2. Tu menú}

Como \manualRol{} el menú lateral te muestra seis opciones:

\begin{center}
\begin{tabularx}{0.9\linewidth}{@{}l >{\raggedright\arraybackslash}X@{}}
\toprule
\textbf{Opción} & \textbf{Para qué te sirve} \\
\midrule
\menu{Dashboard}          & Resumen del turno: qué hay pendiente y qué está en proceso. \\
\menu{Órdenes}            & Consultar las órdenes de producción y entrar a dispensar. \\
\menu{Dispensado}         & \textbf{Tu pantalla principal.} La estación de pesaje. \\
\menu{Recetas}            & Consultar fórmulas. \textbf{No puedes crearlas ni editarlas.} \\
\menu{Inventario}         & Consultar materiales y lotes. \textbf{Sólo consulta.} \\
\menu{Códigos de Barras}  & Imprimir o ver en el celular los códigos de lotes y materiales. \\
\bottomrule
\end{tabularx}
\end{center}

\vspace{6pt}

No verás \menu{Auditoría} ni \menu{Usuarios}: tu rol no tiene acceso a esas pantallas.
Si escribes la dirección a mano, el sistema te regresa al Dashboard.

\subseccion{Lo que un operario NO puede hacer}
\begin{itemize}
    \item Crear órdenes de producción, recetas o materiales.
    \item Aprobar, rechazar o retener lotes.
    \item Firmar una fase (ni la tuya ni la de nadie).
    \item Corregir o borrar un pesaje ya confirmado.
    \item Cancelar una orden.
\end{itemize}

% =====================================================================
\newpage
\seccion{3. Dashboard}

Al entrar verás el resumen del día.

\captura{03_dashboard.png}{Dashboard: tarjetas de conteo, procesos activos, últimas órdenes y actividad reciente.}

\vspace{6pt}

Las tarjetas de arriba son enlaces: al hacer clic te llevan a la pantalla correspondiente.

\begin{itemize}
    \item \textbf{Órdenes Pendientes} y \textbf{En Proceso} $\rightarrow$ \menu{Órdenes}.
    \item \textbf{Dispensados Hoy} $\rightarrow$ \menu{Dispensado}.
    \item \textbf{Lotes en Cuarentena} $\rightarrow$ \menu{Inventario}.
\end{itemize}

El panel \textbf{Actividad Reciente} aparece vacío para tu rol: el detalle de la
bitácora sólo lo ven Supervisión, Calidad, Auditoría, Desarrollo y el Administrador.

% =====================================================================
\newpage
\seccion{4. Estación de dispensado}

\subseccion{4.1 Elegir la orden}

Entra a \menu{Dispensado}. Las órdenes se agrupan en tres bloques:
\textbf{En Proceso}, \textbf{Pendientes} y \textbf{Completadas} (este último viene
cerrado; se abre con la flecha).

\captura{05_dispensado_lista.png}{Lista de órdenes. Cada tarjeta muestra el avance por fases e ingredientes.}

\vspace{6pt}

Cada tarjeta muestra el número de orden, la receta, el lote del producto final, el
progreso de fases firmadas y el progreso de ingredientes. Al pie dice \boton{Iniciar}
si nadie la ha tocado o \boton{Continuar} si ya lleva avance. Haz clic en la tarjeta.

\subseccion{4.2 El recorrido, paso a paso}

Para \textbf{cada ingrediente} el recorrido es siempre el mismo:

\vspace{6pt}
\begin{center}
\begin{tikzpicture}[
    node distance=0.4cm and 1.2cm,
    every node/.style={font=\footnotesize},
    paso/.style={draw=pharmablue, fill=pharmablue!8, rounded corners=6pt,
        minimum width=2.4cm, minimum height=1cm, align=center, text width=2.3cm},
    flecha/.style={-{Stealth[length=6pt]}, thick, color=pharmablue}
]
\node[paso] (scan) {1. Escanear\\el lote};
\node[paso, right=of scan] (pesar) {2. Teclear\\el peso};
\node[paso, right=of pesar] (confirmar) {3. Confirmar};
\node[paso, right=of confirmar] (firma) {4. Firma del\\supervisor\\(al cerrar la fase)};
\draw[flecha] (scan) -- (pesar);
\draw[flecha] (pesar) -- (confirmar);
\draw[flecha] (confirmar) -- (firma);
\end{tikzpicture}
\end{center}
\vspace{6pt}

\begin{cajaOjo}
\textbf{No puedes saltarte pasos ni cambiar el orden.} El sistema te coloca siempre
en el primer ingrediente pendiente de la primera fase sin firmar. Si intentas otro,
responde \cd{No puedes saltarte pasos ni cambiar el orden de la receta.}
\end{cajaOjo}

\newpage
\subseccion{4.3 La pantalla de la estación}

\captura{06_dispensado_paso.png}{Estación: avance por fases, datos del paso actual, rango permitido y campo de escaneo.}

\vspace{6pt}

De arriba hacia abajo verás:

\begin{itemize}
    \item \textbf{Encabezado:} \cd{Dispensado: ORD-00004}, la receta, el lote del
          producto final y el avance (\cd{0/2 fases}, \cd{0/4 ingredientes}).
    \item \textbf{Tira de fases:} una casilla por fase. Verde = firmada, azul = la que
          estás trabajando, gris con candado = todavía bloqueada.
    \item \textbf{Fase activa:} su nombre y, si la receta las trae, las
          \textbf{Instrucciones de fase}.
    \item \textbf{Paso actual:} \cd{PASO 1 DE 2}, el nombre del material, su
          \cd{Código:}, las instrucciones del ingrediente y tres recuadros:
          \textbf{Mínimo}, \textbf{Target} y \textbf{Máximo}, con el porcentaje de
          tolerancia debajo.
\end{itemize}

Si el material es peligroso, el recuadro se pinta de rojo y aparece la etiqueta
\textbf{MATERIAL PELIGROSO}.

\begin{cajaPeligro}
Ponte el equipo de protección que indique el procedimiento de la planta
(mascarilla, guantes, lentes) \textbf{antes} de abrir el contenedor.
\end{cajaPeligro}

\newpage
\subseccion{4.4 Paso 1: escanear el lote}

Dispara la pistola sobre la etiqueta del contenedor, apuntando al \textbf{número de
lote}. La pistola escribe en el campo y manda el Enter sola.

\captura{07_dispensado_scan.png}{Campo de escaneo con el número de lote capturado, listo para Enter.}

\vspace{4pt}

Si no tienes pistola: escribe el código a mano y presiona \boton{Enter}.
También hay un botón de cámara a la derecha del campo para escanear con la cámara
del equipo o del celular.

Al mandar el código, el servidor revisa cuatro cosas:

\begin{enumerate}
    \item Que el lote sea \textbf{de ese material} y no de otro.
    \item Que esté \aprobado{} (un lote en \cuarentena{} o \rechazado{} no sirve).
    \item Que \textbf{no esté caducado}.
    \item Que le quede \textbf{cantidad suficiente} para este pesaje.
\end{enumerate}

Si algo falla, aparece un mensaje rojo bajo el campo y \textbf{no avanzas}:

\captura{29_lote_rechazado.png}{Lote de otro material: el sistema lo rechaza y te deja en el mismo paso.}

\vspace{4pt}

\begin{center}
\begin{tabularx}{0.95\linewidth}{@{}p{0.42\linewidth} >{\raggedright\arraybackslash}X@{}}
\toprule
\textbf{Mensaje} & \textbf{Qué hacer} \\
\midrule
\cd{Material incorrecto...}          & Tomaste el contenedor equivocado. Busca el que dice el paso. \\
\cd{Lote en estado CUARENTENA...}    & Ese lote no está liberado. Avisa a Calidad o a tu supervisor. \\
\cd{Lote CADUCADO (...)}             & Ese lote ya no se puede usar: deja de servir a las 00:00 del día que marca su caducidad. Sácalo de la línea y avisa a Almacén. \\
\cd{Cantidad insuficiente...}        & Ese contenedor ya no alcanza. Pide otro lote a Almacén. \\
\cd{Lote no encontrado...}           & El código se leyó mal o no está dado de alta. Reintenta o avisa. \\
\bottomrule
\end{tabularx}
\end{center}

\begin{cajaOjo}
\textbf{Nunca intentes forzar el dispensado.} Estos bloqueos son del servidor,
no de la pantalla: no se saltan recargando ni escribiendo la dirección a mano.
\end{cajaOjo}

\newpage
\subseccion{4.5 Paso 2: teclear el peso}

Si el lote es correcto, el mensaje de escaneo desaparece y la pantalla cambia sola al
\textbf{Registro de Peso}. Esa es la señal de que el lote quedó validado.

\captura{08_dispensado_lote_ok.png}{Lote aceptado: se abre el Registro de Peso con el semáforo en \cd{Esperando peso}.}

\begin{cajaNota}
\textbf{La báscula no está conectada al sistema.} Pesas en tu balanza y
\textbf{tecleas} el número que lees en la pantalla de la estación. Verifica el
número dos veces antes de confirmar: el sistema no tiene forma de comprobarlo.
\end{cajaNota}

Conforme escribes, el semáforo reacciona:

\begin{center}
\begin{tabularx}{0.9\linewidth}{@{}c l >{\raggedright\arraybackslash}X@{}}
\toprule
& \textbf{Etiqueta en pantalla} & \textbf{Significado} \\
\midrule
\semaverde{}    & \cd{OK - DENTRO DE TOLERANCIA}   & Peso correcto. Botón verde \boton{Confirmar Peso}. \\[4pt]
\semaamarillo{} & \cd{PRECAUCIÓN - Cerca del límite} & Estás dentro del rango, pero pegado a la orilla. Botón amarillo \boton{Confirmar con Precaución}. Se puede confirmar; mejor ajusta. \\[4pt]
\semarojo{}     & \cd{BAJO} o \cd{ALTO - Fuera de tolerancia} & Fuera del rango. Botón gris \textbf{BLOQUEADO}. \\
\bottomrule
\end{tabularx}
\end{center}

\captura{09_peso_verde.png}{Peso dentro de tolerancia: semáforo verde y botón \boton{Confirmar Peso} habilitado.}

\captura{10_peso_rojo.png}{Peso por debajo del mínimo: aviso \cd{FUERA DE TOLERANCIA - NO PERMITIDO} y botón \cd{BLOQUEADO - Ajuste el peso dentro de tolerancia}.}

\vspace{6pt}

El aviso rojo te dice exactamente qué hacer:
\cd{El peso es menor al mínimo permitido (...). Agregue más material.} o bien
\cd{El peso excede el máximo permitido (...). Retire material.}

\captura{11_peso_amarillo.png}{Peso dentro del rango pero cerca del límite: semáforo ámbar.}

\begin{cajaOjo}
\textbf{Un peso fuera de tolerancia no se registra.} No queda a medias ni "pendiente
de autorización": simplemente no se envía. Ajusta el material hasta entrar al rango.
Si la receta no te da el rango que necesitas, llama a tu supervisor: el cambio se
hace en la receta, no en la estación.
\end{cajaOjo}

\newpage
\subseccion{4.6 Paso 3: confirmar}

Al pulsar \boton{Confirmar Peso} el servidor, en una sola operación:

\begin{enumerate}
    \item Vuelve a verificar la tolerancia (aunque la pantalla ya la haya validado).
    \item Guarda el peso real con tu usuario y la hora.
    \item Descuenta esa cantidad del lote.
    \item Avanza al siguiente ingrediente.
\end{enumerate}

\begin{cajaTip}[Doble clic]
Si das doble clic no pasa nada malo: el pesaje se registra \textbf{una sola vez} y
el inventario se descuenta \textbf{una sola vez}. El botón se pone en
\cd{Registrando...} mientras trabaja.
\end{cajaTip}

Un lote nunca queda en negativo. Si al descontar llega a cero, el sistema lo marca
\agotado{} solo, y ya no aparecerá para escanear.

\subseccion{4.7 Cierre de fase: la firma}

Cuando terminas el último ingrediente de una fase, la pantalla cambia sola al
formulario de firma. \textbf{La fase siguiente queda bloqueada hasta que se firme.}

\captura{22_firma_fase.png}{Firma de fase: el supervisor escribe su propio correo y su propia contraseña.}

\begin{enumerate}
    \item Llama a tu supervisor (o a Calidad) a la estación.
    \item \textbf{Él} escribe su correo y su contraseña. Tú no tocas el teclado.
    \item Pulsa \boton{Firmar y Aprobar Fase}.
\end{enumerate}

\begin{cajaOjo}
Un \manualRol{} \textbf{no puede firmar}. Sólo valen las contraseñas de
\textbf{Supervisor}, \textbf{Calidad} o \textbf{Administrador}. Si intentas con la tuya,
el sistema responde \cd{Firma rechazada: credenciales inválidas o el usuario no es
Supervisor, Calidad ni Admin.} y ese rechazo queda anotado en la bitácora.
A los 5 rechazos, ese correo queda bloqueado 15 minutos.
\end{cajaOjo}

Si estás a medio recorrido y sales de la pantalla, no se pierde nada: al volver, la
estación te deja justo en el paso que seguía.

\newpage
\subseccion{4.8 Orden terminada}

Cuando se firma la última fase, la orden pasa a \dispensado{} y la estación muestra
la pantalla de cierre.

\captura{23_orden_completada.png}{Todas las fases dispensadas y firmadas.}

\vspace{6pt}

A partir de ahí la orden ya \textbf{no admite más pesajes}. Si intentas entrar de
nuevo, el sistema responde \cd{La orden está DISPENSADO: ya no admite pesajes.}

% =====================================================================
\newpage
\seccion{5. Consultas que puedes hacer}

\subseccion{5.1 Órdenes}

En \menu{Órdenes} ves la tabla completa de producción.

\captura{04_ordenes.png}{Órdenes de producción con su estado y la acción disponible.}

\vspace{6pt}

La última columna dice \boton{Dispensar} (te lleva a la estación) cuando la orden
está \pendiente{} o \enproceso; en cualquier otro caso dice \boton{Ver detalle} y
abre la vista de consulta. \textbf{No hay botón de Nueva Orden para tu rol}, ni
forma de cancelar una orden desde la pantalla.

\subseccion{5.2 Inventario}

\captura{14_inventario_materiales.png}{Materiales con su stock. Los que están por debajo del mínimo salen resaltados en rojo.}

\vspace{6pt}

En la pestaña \menu{Lotes} ves el estado de cada lote. Para tu rol la tabla es de
\textbf{sólo lectura}: no aparece la columna de acciones ni la pestaña de Recepción.

\captura{15_inventario_lotes.png}{Lotes con cantidad restante, proveedor, caducidad y estado.}

\vspace{4pt}

\begin{itemize}
    \item \aprobado{}: liberado por Calidad. \textbf{Es el único que puedes usar.}
    \item \cuarentena{}: recién recibido o retenido. No se puede usar.
    \item \rechazado{}: descartado por Calidad. No se puede usar.
    \item \agotado{}: se acabó. Lo marca el sistema al llegar a cero.
\end{itemize}

\newpage
\subseccion{5.3 Recetas}

\captura{12_recetas.png}{Recetas con versión, rendimiento, número de fases e ingredientes.}

\vspace{6pt}

Sólo consulta. Las recetas las cargan Desarrollo, Supervisión o el Administrador, y
\textbf{no se editan}: si una fórmula cambia, se da de alta otra.

\subseccion{5.4 Códigos de barras}

\menu{Códigos de Barras} genera las etiquetas de lotes y materiales en pantalla.
Sirve para dos cosas: imprimirlas y pegarlas en los contenedores, o abrir la página
en el celular y escanear desde ahí con la cámara.

% =====================================================================
\newpage
\seccion{6. Qué sí puedes romper}

Cosas que dependen de ti y que el sistema \textbf{no} puede corregir:

\begin{center}
\begin{tabularx}{\linewidth}{@{}p{0.34\linewidth} >{\raggedright\arraybackslash}X@{}}
\toprule
\textbf{Riesgo} & \textbf{Cómo evitarlo} \\
\midrule
\textbf{Teclear mal el peso} & La báscula no está conectada: el número que escribes es el que queda registrado como peso real. Léelo dos veces. Un pesaje confirmado \textbf{no se puede corregir ni borrar} desde la pantalla. \\
\textbf{Escanear un contenedor y pesar de otro} & El sistema valida el \textit{código}, no el polvo. Escanea el contenedor del que realmente vas a tomar el material. \\
\textbf{Prestar tu sesión} & Todo queda a tu nombre. Si te alejas de la estación, cierra sesión con el botón junto a tu nombre, abajo del menú. \\
\textbf{Que el supervisor firme sin ver} & La firma certifica lo que hiciste. Muéstrale los contenedores y el resumen de la fase antes de que teclee. \\
\textbf{Dejar una fase sin firmar} & La orden se queda atorada y nadie puede continuarla. Consigue la firma antes de terminar tu turno. \\
\bottomrule
\end{tabularx}
\end{center}

% =====================================================================
\newpage
\seccion{7. Preguntas frecuentes}

\subseccion{Escaneé el lote equivocado, ¿ya se registró algo?}
No. Nada se guarda hasta que confirmas el peso. Escanea el lote correcto y sigue.

\subseccion{Se me pasó la mano y el peso quedó alto}
Retira material hasta entrar al rango. El botón se desbloquea solo cuando el
semáforo deja de estar rojo.

\subseccion{Confirmé un peso mal. ¿Lo puedo corregir?}
No desde la pantalla. Avisa de inmediato a tu supervisor y a Calidad: la corrección
se maneja fuera del sistema, y el pesaje original queda en la bitácora.

\subseccion{¿Puedo pausar y continuar después?}
Sí. El avance se guarda en el servidor. Al volver a la orden retomas en el paso
pendiente, aunque sea en otra computadora.

\subseccion{¿Quién puede firmar mi fase?}
Supervisor, Calidad o Administrador, con su propia contraseña. Nunca tú.

\subseccion{Me bloquearon la cuenta}
Fueron 5 intentos fallidos con tu correo. Espera 15 minutos. No hay botón de
desbloqueo ni de "olvidé mi contraseña": si de verdad la perdiste, el administrador
tiene que generarte una nueva desde el servidor.

\subseccion{¿Dónde veo lo que ya pesé en esta fase?}
Debajo del campo de peso, en \textbf{Ingredientes de esta fase}: los ya dispensados
salen con palomita y su peso real; los pendientes, con su target en gris.
